\section*{S3. Per-family results}

Table~S1 breaks the XGBoost results of the main text down by attack family (all families that appear in the 1,000-flow cohort; families with a handful of flows, such as one CIC-IDS2017 botnet flow, give no reliable rate). On CSE-CIC-IDS2018 the families in the cohort are evaded in 98.7\%--100.0\% of the flows (DoS 390 flows, DDoS 375, brute force 200, botnet 35). On CIC-IDS2017 DoS (422 flows) is evaded in 98.1\% and the other families in 100.0\%. On UNSW-NB15 more than half of the cohort are exploits, which are evaded in only 44.6\% of 504 flows and so set the aggregate figure; generic attacks (50.0\% on 22 flows) and DoS (51.2\% on 80) are similar, whereas fuzzers (79.5\% on 219) and reconnaissance (69.2\% on 143) are easier.

\begin{table}[!htbp]
\centering
\caption*{Table S1. FFE against XGBoost by attack family.}
\begin{tabular}{@{}llccc@{}}
\toprule
\textbf{Dataset} & \textbf{Attack family} & \textbf{$n$} & \textbf{Real ASR (\%)} & \textbf{Queries} \\
\midrule
CSE-CIC-IDS2018 & DoS & 390 & 100.0 & 7.6 \\
 & DDoS & 375 & 98.7 & 29.2 \\
 & BruteForce & 200 & 100.0 & 9.9 \\
 & Botnet & 35 & 100.0 & 18.4 \\
\midrule
CIC-IDS2017 & DoS & 422 & 98.1 & 8.0 \\
 & PortScan & 322 & 100.0 & 29.4 \\
 & DDoS & 235 & 100.0 & 5.0 \\
 & BruteForce & 20 & 100.0 & 28.0 \\
 & Botnet & 1 & 100.0 & 9.0 \\
\midrule
UNSW-NB15 & Exploits & 504 & 44.6 & 116.7 \\
 & Fuzzers & 219 & 79.5 & 56.4 \\
 & Reconnaissance & 143 & 69.2 & 69.4 \\
 & DoS & 80 & 51.2 & 104.7 \\
 & Generic & 22 & 50.0 & 117.2 \\
 & Shellcode & 19 & 100.0 & 13.8 \\
 & Backdoor & 6 & 66.7 & 103.5 \\
 & Worms & 6 & 33.3 & 131.0 \\
 & Analysis & 1 & 100.0 & 9.0 \\
\bottomrule
\end{tabular}

\end{table}

\section*{S4. Ablations}

Table~S2 compares FFE with and without tree-split boundary snapping (TSBS) on the full 1,000-flow cohorts of the three corpora. Table~S3 reports the other ablations on the CUDA XGBoost target of CSE-CIC-IDS2018 with a balanced cohort of 283 detected malicious flows (up to 50 per attack family): Panel A varies the step budget $M$, Panel B isolates generator subsets, and Panel C compares operational cost regimes. The TSBS relative snapping offset $\epsilon_{\text{rel}}$ was swept over $\{0, 10^{-6}, 10^{-4}, 10^{-2}, 1\}$.

\begin{table}[!htbp]
\centering
\caption*{Table S2. FFE with and without tree-split boundary snapping (XGBoost, Real ASR with Wilson 95\% interval).}
\begin{tabular}{@{}lcccc@{}}
\toprule
\textbf{Dataset} & \textbf{Variant} & \textbf{Real ASR (\%)} & \textbf{Queries} & \textbf{ms/flow} \\
\midrule
CSE-CIC-IDS2018 & FFE & 99.5 [98.8, 99.8] & 16.5 & 15.8 \\
 & FFE without TSBS & 99.5 [98.8, 99.8] & 11.6 & 10.6 \\
CIC-IDS2017 & FFE & 99.2 [98.4, 99.6] & 14.6 & 10.7 \\
 & FFE without TSBS & 97.5 [96.3, 98.3] & 13.5 & 9.6 \\
UNSW-NB15 & FFE & 57.6 [54.5, 60.6] & 93.7 & 89.1 \\
 & FFE without TSBS & 54.6 [51.5, 57.7] & 51.5 & 36.0 \\
\bottomrule
\end{tabular}

\end{table}

\textbf{Step budget.} ASR rises from 42.8\% at $M=1$ (5.0 queries) through 67.1\% at $M=3$ and 84.1\% at $M=5$ to 98.9\% at $M=10$ and $M=20$ (14.4 queries at $M=20$). Physical feasibility is 100.0\% for every budget.

\textbf{Generators.} No single generator suffices. Temporal dilation alone evades 66.1\% of the flows but needs 32.6 queries and adds 40.7\,s of flow duration on average; window morphing alone evades 46.3\% with 2.5 queries and no added delay; chaff alone 32.5\%; padding alone 10.2\%. Pairs are better than singles (time plus chaff 78.1\%, time plus padding 69.3\%, padding plus chaff 33.6\% with no added delay), and the complete cone reaches 98.9\% with 14.4 queries and 0.18\,s of mean delay. The full cone needs fewer queries than the time-based subsets, possibly because the cheap window and chaff moves are available when delay alone would stall.

\textbf{Cost regimes and snapping offset.} These do not discriminate on this corpus. All four cost regimes reach 98.9\% ASR with 14.1--14.4 queries, and the snapping offset leaves ASR at 98.9\% and the queries between 14.2 and 14.4. The regimes do change the mean added delay (181\,ms under regime A, 133\,ms under the latency-sensitive regime B), so FFE responds to the operator's cost model, but because the corpus is saturated we have no evidence on how the regime affects success; a harder target such as UNSW-NB15 would be needed for that.

\begin{table}[!htbp]
\centering
\caption*{Table S3. Ablation of FFE on CUDA XGBoost, CSE-CIC-IDS2018 ($N=283$). Delay is the mean added flow duration.}
\resizebox{\columnwidth}{!}{%
\begin{tabular}{@{}lcccc@{}}
\toprule
\textbf{Configuration} & \textbf{ASR (\%)} & \textbf{PFR (\%)} & \textbf{Queries} & \textbf{Added delay (s)} \\
\midrule
\multicolumn{5}{@{}l}{\textit{Panel A: step budget $M$}} \\
$M=1$ & 42.8 & 100.0 & 5.0 & 0.08 \\
$M=2$ & 63.3 & 100.0 & 8.3 & 0.14 \\
$M=3$ & 67.1 & 100.0 & 9.5 & 0.15 \\
$M=5$ & 84.1 & 100.0 & 12.2 & 0.16 \\
$M=10$ & 98.9 & 100.0 & 13.9 & 0.18 \\
$M=20$ & 98.9 & 100.0 & 14.4 & 0.18 \\
$M=50$ & 98.9 & 100.0 & 16.0 & 0.18 \\
\midrule
\multicolumn{5}{@{}l}{\textit{Panel B: generator subset ($M=20$)}} \\
Temporal Dilation Only & 66.1 & 100.0 & 32.6 & 40.72 \\
Payload Padding Only & 10.2 & 100.0 & 5.1 & 0.00 \\
Chaff Injection Only & 32.5 & 100.0 & 9.7 & 0.00 \\
Window Morphing Only & 46.3 & 100.0 & 2.5 & 0.00 \\
Time + Padding & 69.3 & 100.0 & 39.6 & 26.81 \\
Time + Chaff & 78.1 & 100.0 & 38.4 & 22.05 \\
Padding + Chaff & 33.6 & 100.0 & 14.2 & 0.00 \\
Complete Lie Cone & 98.9 & 100.0 & 14.4 & 0.18 \\
\midrule
\multicolumn{5}{@{}l}{\textit{Panel C: operational cost regime ($M=20$)}} \\
A (Default / Balanced) & 98.9 & 100.0 & 14.4 & 0.18 \\
B (Latency-Sensitive) & 98.9 & 100.0 & 14.1 & 0.13 \\
C (Unconstrained Delay) & 98.9 & 100.0 & 14.4 & 0.18 \\
D (Uniform / Flat) & 98.9 & 100.0 & 14.3 & 0.14 \\
\bottomrule
\end{tabular}
%
}
\end{table}

\section*{S5. Cross-architecture transferability}

Flows crafted by FFE against a source ensemble were scored by the three other architectures (Table~S4). The cohort is the 221 CSE-CIC-IDS2018 malicious flows (up to 50 per family) that XGBoost, LightGBM, CatBoost and Random Forest all detect when clean, so a transferred evasion is not the target's own false negative; all four models are trained on the same training partition. Flows crafted against XGBoost evade LightGBM in 71.5\% of the cases, Random Forest in 55.7\% and CatBoost in 24.0\%. LightGBM flows evade XGBoost in 75.6\%, Random Forest in 63.8\% and CatBoost in 24.0\%; CatBoost flows evade XGBoost in 70.1\%, LightGBM in 77.4\% and Random Forest in 77.4\%; Random Forest flows evade XGBoost in 78.7\%, LightGBM in 71.0\% and CatBoost in 24.0\%. Excluding CatBoost as target, 55.7\%--78.7\% of the transferred flows evade. CatBoost is the hardest target for the three other sources, whose flows each evade it in 53 of 221 cases (we did not check whether these are the same flows); attacks against CatBoost itself succeed in only 69.7\% of the cases. We did not test why CatBoost differs; its symmetric oblivious trees place thresholds differently from the leaf-wise and depth-wise ensembles.

\begin{table}[!htbp]
\centering
\caption*{Table S4. Transferability of FFE flows (attack success rate in \%, $N=221$; diagonal: attack against the source model itself).}
\resizebox{\columnwidth}{!}{%
\begin{tabular}{@{}lcccc@{}}
\toprule
\textbf{Source $\backslash$ Target} & \textbf{XGBoost (CUDA)} & \textbf{LightGBM} & \textbf{CatBoost} & \textbf{Random Forest} \\
\midrule
\textbf{XGBoost (CUDA)} & \textbf{100.0} & 69.0 & 37.7 & 81.3 \\
\textbf{LightGBM} & 82.3 & \textbf{100.0} & 34.0 & 83.0 \\
\textbf{CatBoost} & 54.3 & 66.7 & \textbf{99.3} & 63.7 \\
\textbf{Random Forest} & 79.0 & 64.0 & 37.0 & \textbf{100.0} \\
\bottomrule
\end{tabular}
%
}
\end{table}

\section*{S6. Corpus statistics}

Table~S5 gives the record counts of the three main corpora and of the three additional corpora of Section~S7. Valid records pass the validity filter; Kept = Valid $-$ Duplicates, where duplicates are records with identical features and label after projection onto the flow manifold. Records with identical features but different labels are kept (149,526, 22 and 1,040).

\begin{table}[!htbp]
\centering
\caption*{Table S5. Full-scale benchmark corpora (the last three are those of Section~S7).}
\footnotesize
\begin{tabular}{@{}lrrrrrrr@{}}
\toprule
\textbf{Dataset} & \textbf{Raw rows} & \textbf{Valid} & \textbf{Duplicates} & \textbf{Kept} & \textbf{Attacks} & \textbf{Train} & \textbf{Test} \\
\midrule
CSE-CIC-IDS2018 & 8,284,254 & 8,153,708 & 3,185,319 & 4,968,389 & 594,809 & 3,477,872 & 993,679 \\
CIC-IDS2017 & 2,099,971 & 2,085,851 & 426,982 & 1,658,869 & 396,403 & 1,161,208 & 331,774 \\
UNSW-NB15 & 257,673 & 215,689 & 92,181 & 123,508 & 55,559 & 86,456 & 24,703 \\
\bottomrule
\end{tabular}

\end{table}

\section*{S7. Three additional public corpora}

We ran the same pipeline (cleaning, projection, deduplication, 70/10/20 split, XGBoost and LightGBM with 60 trees of depth 6, a cohort of up to 1,000 detected flows, one seed) on three further public corpora. The builders are in \texttt{src/datasets/full\_data.py}; each needed dataset-specific decisions that we list because they bound what the results mean.

\textbf{InSDN} (software-defined networking; three CSV files, 343,889 rows; CC BY-NC 4.0). The exporter wrote $-1$ as the TCP window of every row, so the window is the constant 0. Almost all DDoS flows (about 122,000) carry protocol 0, which is neither TCP nor UDP, and are removed by the validity filter, so the DDoS family is absent. After cleaning, 130,076 flows remain (58,935 benign, 37,822 DoS, 31,770 probe, 1,210 brute force, 164 botnet, 158 web, 17 U2R).

\textbf{USB-IDS-1} (17 files: four DoS tools times four server defenses, plus regular traffic; 4,815,203 rows). After cleaning, 3,874,467 flows remain, of which 92\% are attacks (Hulk 84\% of all flows), because the corpus is built around DoS attacks. Both XGBoost and LightGBM reach 100.00\% accuracy.

\textbf{CIDDS-001} (unidirectional NetFlow, OpenStack and external server, 31,959,174 rows). NetFlow records have no backward direction or TCP window, so backward packets and bytes are 0 and the window is 0 for every flow; payload bytes are the byte count minus 20 bytes of IP header and 20 (TCP) or 8 (UDP) bytes of transport header per packet. We keep the class \texttt{normal} as benign and \texttt{attacker} as attack; the \texttt{victim} flows (responses of the attacked host), \texttt{suspicious} and \texttt{unknown} are ambiguous and dropped. After cleaning, 7,084,648 flows remain, of which only 5,560 are attacks (0.08\%: port scan 3,171, brute force 1,329, DoS 1,000, ping scan 60). XGBoost reaches F1 94.4\% and LightGBM only 66.7\%; 819 of the 1,112 test attacks are detected by both models and form the cohort.

Table~S6 gives Real ASR and mean queries against XGBoost. FFE is not uniformly the best method on these corpora.
\begin{itemize}
\item \textbf{InSDN.} FFE reaches 81.0\% (79.9\% for CAPGD, 85.5\% for PSO on 200 flows). FFE beats CAPGD (11 flows against none, Holm-adjusted $p=0.006$) but is not distinguishable from PSO on the 200 flows both attacked (82.5\% against 85.5\%; 2 against 8 discordant flows, adjusted $p=0.44$). Here snapping matters: FFE without TSBS reaches only 46.7\%.
\item \textbf{USB-IDS-1.} FFE and PSO reach 0.0\%, whereas the unconstrained attack, ZOO, MOEVA, CAPGD and the generator reach 95.3\%--99.2\%. The target's score is the same for every attack flow we looked at (margin 8.57 for the first 5,000 test attacks), so score probes give FFE no direction. CAPGD's successful flows change four things at once (forward packets, duration, window, header length; in one example the window fell from 29,200 to 21,368), while FFE's window action only moves toward a fixed target. We did not isolate which of these limits FFE.
\item \textbf{CIDDS-001.} FFE reaches 100.0\% with 6.2 queries; PSO also reaches 100.0\% (no discordant flow among 200), and the generator 99.1\% (7 flows against none, adjusted $p=0.08$). The attacks are few, the corpus is unidirectional, and nearly every method that repairs its output succeeds, so this corpus separates the methods little.
\end{itemize}
On these three corpora FFE therefore wins or ties on one, is statistically indistinguishable from the best baseline on one, and fails on one. The claim in the main text that FFE has the highest Real ASR applies to the three main corpora only.

\begin{table}[!htbp]
\centering
\caption*{Table S6. Real ASR (\%, Wilson 95\% interval) and mean target queries against XGBoost on the additional corpora; MOEVA and PSO attack the first 200 flows of the cohort.}
\resizebox{\textwidth}{!}{%
\begin{tabular}{@{}lcccccc@{}}
\toprule
\textbf{Method} & \multicolumn{2}{c}{\textbf{InSDN}} & \multicolumn{2}{c}{\textbf{USB-IDS-1}} & \multicolumn{2}{c}{\textbf{CIDDS-001}} \\
 & Real ASR [95\% CI] & Queries & Real ASR [95\% CI] & Queries & Real ASR [95\% CI] & Queries \\
\midrule
XGBoost acc.\ / F1 (\%) & \multicolumn{2}{c}{99.93 / 99.93 (cohort 1000)} & \multicolumn{2}{c}{100.00 / 100.00 (cohort 1000)} & \multicolumn{2}{c}{99.99 / 94.42 (cohort 819)} \\
\midrule
Unconstrained feature space & 0.0 [0.0, 0.4] & 47.0 & 99.2 [98.4, 99.6] & 27.2 & 0.6 [0.3, 1.4] & 27.9 \\
ZOO-NIDS & 0.0 [0.0, 0.4] & 107.7 & 95.3 [93.8, 96.4] & 39.5 & 0.0 [0.0, 0.5] & 126.9 \\
HopSkipJump-NIDS & 0.0 [0.0, 0.4] & 47.0 & 0.0 [0.0, 0.4] & 47.0 & 0.0 [0.0, 0.5] & 47.0 \\
MOEVA (search stage of CAA) & 0.0 [0.0, 1.9] & 501.6 & 98.5 [95.7, 99.5] & 31.8 & 7.0 [4.2, 11.4] & 441.6 \\
PSO (Traffic Manipulator, projected) & 85.5 [80.0, 89.7] & 72.0 & 0.0 [0.0, 1.9] & 377.0 & 100.0 [98.1, 100.0] & 17.0 \\
CAPGD (surrogate) & 79.9 [77.3, 82.3] & 5.6 & 99.1 [98.3, 99.5] & 1.2 & 86.9 [84.5, 89.1] & 3.5 \\
NIDSGAN-style & 38.7 [35.7, 41.8] & 3.5 & 99.1 [98.3, 99.5] & 1.0 & 99.1 [98.2, 99.6] & 1.0 \\
\textbf{FFE (ours)} & 81.0 [78.5, 83.3] & 47.4 & 0.0 [0.0, 0.4] & 181.0 & 100.0 [99.5, 100.0] & 6.2 \\
FFE without TSBS & 46.7 [43.6, 49.8] & 59.9 & 0.0 [0.0, 0.4] & 101.0 & 100.0 [99.5, 100.0] & 6.2 \\
\bottomrule
\end{tabular}
%
}
\end{table}
